\section{Introduction \& Motivation}

\begin{frame}[allowframebreaks]{Why Audio in NLP?}
  \large{Speech is the \textbf{most natural form of communication}.}

  \vspace{0.5em}

  \normalsize
  \textbf{What:} Where speech interfaces are essential.
  \begin{itemize}
    \item Voice assistants (e.g., Amazon Echo, Siri, Google Assistant)
    \item Accessibility: captions, screen readers
    \item Hands-free operation (cars, IoT)
    \item Multilingual interaction \& translation
  \end{itemize}
  \vspace{0.5em}
  \textbf{Why:} Speech enables more intuitive and inclusive interaction.

\framebreak

  \textbf{Why audio needs its own toolkit.}
  \begin{itemize}
    \setlength{\itemsep}{1em}
    \item Audio data is abundant and diverse, encompassing speech, music, and environmental sounds.
    \item Text-oriented methods do not apply directly: audio arrives as a continuous waveform with no tokens, no word boundaries, and thousands of samples per second.
    \item Many NLP systems now start from speech, not just text (e.g., live captions, real-time translation).
    \item The signal carries information text discards: speaker identity, prosody, emotion, and overlapping speakers.
  \end{itemize}
\end{frame}
